\section{Numerical simulations}
\label{sec:sim}

In this section, we check the performance of the classifiers and the validity of
the theoretical results by numerical simulations. We give five experiments. In
Section~\ref{sec:sim-thm4}, we verify Theorem~\ref{thm:4} in the setting~(E). In
Section~\ref{sec:sim-gram}, we verify the convergence of the Gram matrix given in
Theorem~\ref{thm:1}. In Section~\ref{sec:sim-inconsist}, we check the inconsistency
given in Corollary~\ref{cor:1}. In Section~\ref{sec:sim-bias}, we check the bias
term given in Proposition~\ref{prop:1}. In Section~\ref{sec:sim-classifiers}, we
compare the classifiers including the SC-SVM.

Throughout this section, the populations are Gaussian and we take the coordinates
so that the first coordinate is the direction of $\mu$, the next $m_i$ coordinates
are the spiked directions, which are orthogonal to $\mu$, and the remaining
coordinates give the isotropic non-spiked part. That is,
\begin{equation}
\Sigma_i=\diag(\sigma_i^2,\lambda_{i1},\dots,\lambda_{im_i},\sigma_i^2,\dots,\sigma_i^2),
\quad
\lambda_{is}=c_{is}\theta,
\quad
\theta=d,
\label{eq:11.1}
\end{equation}
where $\sigma_i^2$ is determined by $\tr(\Sigma_{i*})=\theta\,c_{i0}$. Since
$\Sigma_i$ is diagonal, we can compute the conditional misclassification rates
exactly for each realization of the training data, so that we do not need a test
sample. We give the standard errors of the averages over the replications.
The code for each experiment is distributed with this manuscript as a separate
Python script (\texttt{spiked\_svm\_sim.py}, \texttt{spiked\_svm\_gram.py},
\texttt{spiked\_svm\_inconsistency.py}, \texttt{spiked\_svm\_bias.py},
\texttt{spiked\_svm\_classifiers.py}), together with the shared utilities in
\texttt{spiked\_svm\_common.py}.

\subsection{Verification of Theorem~\ref{thm:4}}
\label{sec:sim-thm4}

We first consider the setting~(E) with $d=20000$, in which the limiting
misclassification rate is given by~\eqref{eq:7.2} in a closed form. We considered
$c\in\{0.1,0.3,0.5,0.8\}$ and $\delta\in\{0,1,4\}$, and took the average over
$600$ replications. We computed the limiting value~\eqref{eq:7.2} by the Monte
Carlo integration with $2\times 10^5$ points. The implementation is given in
\texttt{spiked\_svm\_sim.py}. We give the results in
Table~\ref{tab:2}.

\begin{table}[ht]
\centering
\caption{The limiting misclassification rate~\eqref{eq:7.2} and the simulated
$E\{e(1)\}$ for $d=20000$ in the setting~(E).}
\label{tab:2}
\begin{tabular}{@{}rrrrr@{}}
\toprule
$c$ & $\delta$ & limit~\eqref{eq:7.2} & simulation & difference \\
\midrule
0.1 & 0.0 & 0.5009 & 0.4993 & 0.0016 \\
0.1 & 1.0 & 0.0112 & 0.0122 & 0.0010 \\
0.1 & 4.0 & 0.0000 & 0.0000 & 0.0000 \\
0.3 & 0.0 & 0.5009 & 0.4998 & 0.0011 \\
0.3 & 1.0 & 0.1093 & 0.1091 & 0.0002 \\
0.3 & 4.0 & 0.0039 & 0.0047 & 0.0008 \\
0.5 & 0.0 & 0.5009 & 0.4997 & 0.0012 \\
0.5 & 1.0 & 0.1833 & 0.1832 & 0.0001 \\
0.5 & 4.0 & 0.0215 & 0.0225 & 0.0010 \\
0.8 & 0.0 & 0.5009 & 0.4997 & 0.0012 \\
0.8 & 1.0 & 0.2517 & 0.2519 & 0.0002 \\
0.8 & 4.0 & 0.0598 & 0.0602 & 0.0004 \\
\bottomrule
\end{tabular}
\end{table}

We observe that the differences are smaller than $0.002$ for all the cases, which
are within the Monte Carlo error whose standard error is about $0.006$ for $600$
replications. Hence Theorem~\ref{thm:4} is confirmed numerically. When
$\delta=0$, the misclassification rate is about $0.5$ for all $c$, so that the
classification is completely random. When $\delta$ is fixed, the misclassification
rate increases monotonically as $c$ increases. We emphasize that the spike is
orthogonal to $\mu$ in this setting, so that it carries no information about the
classification. Nevertheless it degrades the performance of the SVM severely.

\subsection{Convergence of the Gram matrix}
\label{sec:sim-gram}

Next, we check Theorem~\ref{thm:1} directly. We considered $n_1=n_2=2$, $m_i=1$,
$c=0.5$ and $\delta=1$, and took $d$ from $10^2$ to $10^6$. In each replication,
we first generated the spiked coordinates $\{z_{ij1}\}$ and fixed them, and then
generated the non-spiked part for the given $d$. We computed the limiting Gram
matrix $G^*$ from~\eqref{eq:4.3} using the same $\{z_{ij1}\}$. Since the spike is
orthogonal to $\mu$, we have $\beta_{is}=0$ and $r_{11}=1$, so that
\[
g^*_{(ij),(kl)}
=\frac{\eps_i\eps_k\delta}{4}+c\,z_{ij1}z_{kl1},
\qquad
G^*=G_0^*+(1-c)I_N.
\]
We measured $\|G-G^*\|_F$ and took the average over $100$ replications. The
implementation is given in \texttt{spiked\_svm\_gram.py}. We give
the results in Table~\ref{tab:gram}.

\begin{table}[ht]
\centering
\caption{The Frobenius norm $\|G-G^*\|_F$ and the ratio of the successive values.}
\label{tab:gram}
\begin{tabular}{@{}rrrr@{}}
\toprule
$d$ & $\|G-G^*\|_F$ & s.e. & ratio to previous \\
\midrule
100 & 0.29956 & 0.00783 & -- \\
1{,}000 & 0.09099 & 0.00245 & 3.292 \\
10{,}000 & 0.03138 & 0.00077 & 2.900 \\
100{,}000 & 0.01039 & 0.00027 & 3.021 \\
1{,}000{,}000 & 0.00299 & 0.00007 & 3.471 \\
\bottomrule
\end{tabular}
\end{table}

We observe that $\|G-G^*\|_F$ tends to zero. Since the dimension is multiplied by
$10$ in each row, the ratio should be $\sqrt{10}\approx 3.162$ if the error is of
order $d^{-1/2}$. The observed ratios are $3.292$, $2.900$, $3.021$ and $3.471$,
which are close to $\sqrt{10}$. Hence Theorem~\ref{thm:1} is confirmed, and
moreover the convergence rate is of order $d^{-1/2}$ in this setting. We note that
the rate itself is not claimed in Theorem~\ref{thm:1}, which gives the convergence
only.

\subsection{The inconsistency}
\label{sec:sim-inconsist}

We check Corollary~\ref{cor:1}. We considered $n_1=n_2=2$, $m_i=1$ and
$\delta=0.25$, and took $c\in\{0.8,0.5,0.2\}$ together with the non-spiked case in
which $c=10^{-4}$. We took the average of $e(1)+e(2)$ over $200$ replications for
$d\in\{250,1000,4000,16000,64000\}$. The implementation is given in
\texttt{spiked\_svm\_inconsistency.py}. We give the results in
Figure~\ref{fig:inconsist}.

\begin{figure}[ht]
\centering
\includegraphics[width=0.72\textwidth]{fig_inconsistency.pdf}
\caption{Average of $e(1)+e(2)$ for the hard-margin SVM vs.\ $d$
($n_1=n_2=2$, $\delta=0.25$).}
\label{fig:inconsist}
\end{figure}

In the non-spiked case, the misclassification rate tends to zero rapidly as $d$
increases, which is the consistency property of the existing theory. On the other
hand, under the spiked model, the misclassification rate does not decrease at all
even though $d$ is multiplied by $256$. Hence Corollary~\ref{cor:1} is confirmed.

We also observe that the misclassification rate is not monotone in $c$, since the
value for $c=0.5$ is larger than the one for $c=0.8$. We note that the
monotonicity was observed in Table~\ref{tab:2}, in which $n_1=n_2=1$. Hence the
monotonicity does not hold in general when $n_i\ge 2$.

We give a supplementary result in order to clarify the practical meaning of
Corollary~\ref{cor:1}. We took $d=4000$ and $\delta=0.25$, and varied
$n_1=n_2=n$. We give the results in Table~\ref{tab:n-vary}.

\begin{table}[ht]
\centering
\caption{The average of $e(1)+e(2)$ for $d=4{,}000$ and $\delta=0.25$.}
\label{tab:n-vary}
\begin{tabular}{@{}rrrr@{}}
\toprule
$n$ & $c=0.8$ & $c=0.5$ & non-spiked \\
\midrule
2 & 0.2684 & 0.3467 & 0.0000 \\
3 & 0.0771 & 0.1589 & 0.0000 \\
5 & 0.0154 & 0.0424 & 0.0000 \\
10 & 0.0000 & 0.0003 & 0.0000 \\
\bottomrule
\end{tabular}
\end{table}

Although Corollary~\ref{cor:1} states that the limit is positive for any fixed $n$,
the value decreases rapidly as $n$ increases and it is already smaller than
$10^{-3}$ for $n=10$ in this setting. Hence the practical effect of the
inconsistency is restricted to the case that the sample size is very small. We
note that the setting of Table~\ref{tab:n-vary} is favourable in the sense that
$\delta$ is not small. We give a harder setting in
Section~\ref{sec:sim-classifiers}.

\subsection{The bias term}
\label{sec:sim-bias}

We check Proposition~\ref{prop:1}. In the $u$-scale, the midpoint of $\mu_1$ and
$\mu_2$ is the origin, so that the intercept $b$ is exactly the value of the
discriminant function at the midpoint, that is, the deterministic offset. We
measured $b$ and compared it with the predictions given by $\kappa$ and by
$\kappa_*$.

From the outline of the proof of Proposition~\ref{prop:1}, the offset is of the
form $-A\kappa_*/(2\theta)$. Writing $A=\sum_j\alpha_{1j}=\sum_j\alpha_{2j}$,
which holds from the constraint $y^T\alpha=0$, and noting that the dual weights are
equalized within each class, the prediction is
\begin{equation}
b\approx -\frac{A}{2}\cdot\frac{\kappa_*}{\theta}
\quad\text{instead of}\quad
b\approx -\frac{A}{2}\cdot\frac{\kappa}{\theta}.
\label{eq:11.2}
\end{equation}
We note that $A$ is observable, so that~\eqref{eq:11.2} gives a quantitative
prediction without any unknown constant. In order to distinguish the two
predictions, we chose the configurations in which $\kappa$ and $\kappa_*$ disagree
qualitatively. Since $\theta_1=\theta_2=\theta$ in our setting, we have
\[
\frac{\kappa}{\theta}=\frac{1}{n_1}-\frac{1}{n_2},
\qquad
\frac{\kappa_*}{\theta}=\frac{1-c_1}{n_1}-\frac{1-c_2}{n_2},
\]
so that we can make $\kappa=0$ while $\kappa_*\neq 0$ by taking $n_1=n_2$ and
$c_1\neq c_2$, and we can make $\kappa_*=0$ while $\kappa\neq 0$ by a suitable
choice. We took $d=8000$, $\delta=1$ and $300$ replications. The implementation
is given in \texttt{spiked\_svm\_bias.py}. We give the results
in Table~\ref{tab:bias}.

\begin{table}[ht]
\centering
\caption{The measured intercept $b$ and the two predictions in~\eqref{eq:11.2}.}
\label{tab:bias}
\setlength{\tabcolsep}{3pt}
\begin{tabular}{@{}rrrrrrrrrr@{}}
\toprule
$n_1$ & $n_2$ & $c_1$ & $c_2$ & $\kappa/\theta$ & $\kappa_*/\theta$
  & $b$ (s.e.) & $A$ & pred.\ $\kappa$ & pred.\ $\kappa_*$ \\
\midrule
5 & 5 & 0.2 & 0.8 & 0.0000 & 0.1200 & $-0.0899$ (0.0017) & 1.644 & 0.0000 & $-0.0987$ \\
4 & 4 & 0.1 & 0.7 & 0.0000 & 0.1500 & $-0.0921$ (0.0029) & 1.492 & 0.0000 & $-0.1119$ \\
10 & 5 & 0.4 & 0.7 & $-0.1000$ & 0.0000 & $0.0071$ (0.0011) & 1.766 & 0.0883 & 0.0000 \\
8 & 4 & 0.3 & 0.65 & $-0.1250$ & 0.0000 & $0.0103$ (0.0015) & 1.669 & 0.1043 & 0.0000 \\
8 & 4 & 0.5 & 0.5 & $-0.1250$ & $-0.0625$ & $0.0557$ (0.0012) & 1.657 & 0.1035 & 0.0518 \\
\bottomrule
\end{tabular}
\end{table}

We observe the following. In the first two rows, we have $\kappa=0$, so that the
prediction by $\kappa$ is zero. The measured values are $-0.0899$ and $-0.0921$,
which are far from zero compared with the standard errors, while the predictions
by $\kappa_*$ are $-0.0987$ and $-0.1119$. In the third and the fourth rows, we
have $\kappa_*=0$, and the measured values are $0.0071$ and $0.0103$, which are
much smaller than the predictions by $\kappa$, that is, $0.0883$ and $0.1043$. In
the last row, in which both are non-zero, the measured value $0.0557$ is close to
the prediction by $\kappa_*$, that is, $0.0518$, while the prediction by $\kappa$
is $0.1035$, which is about twice as large.

Hence Proposition~\ref{prop:1} is confirmed. The remaining discrepancy, which is
at most about $0.02$ in the first two rows and which is not zero in the third and
the fourth rows, is considered to come from the finite-dimensional effect and from
the fact that the dual weights are not exactly equalized within each class. We
note that the qualitative conclusion is not affected by this discrepancy, because
the two predictions differ in sign or in magnitude by a factor of two.

\subsection{Comparison of the classifiers}
\label{sec:sim-classifiers}

Finally, we compare the classifiers. We consider the following procedures. For
the bias correction, we use the form suggested by the outline of the proof of
Proposition~\ref{prop:1}, that is, we shift the intercept by
\begin{equation}
b\;\longmapsto\;
b+\frac12\bigl(\bar\alpha_1\,\hat c_{10}-\bar\alpha_2\,\hat c_{20}\bigr),
\label{eq:11.3}
\end{equation}
where $\bar\alpha_i$ is the average of the dual weights of class $i$ and
$\hat c_{i0}$ is an estimator of the noise floor.

\begin{itemize}
\item SVM: the hard-margin SVM without any correction.
\item BC-SVM($\kappa$): the correction~\eqref{eq:11.3} with
$\hat c_{i0}=\tr(S_i)/\hat\theta$, where $S_i$ is the sample covariance matrix.
This uses the full trace and corresponds to the existing BC-SVM.
\item BC-SVM($\kappa_*$): the correction~\eqref{eq:11.3} with
$\hat c_{i0}=\{\tr(S_i)-\sum_{s\le m_i}\tilde\lambda_{is}\}/\hat\theta$, where
$\tilde\lambda_{is}$ is the NR estimator of the spiked eigenvalue. This
corresponds to Proposition~\ref{prop:1}.
\item SC-SVM: the procedure given in Definition~\ref{def:2}. We give two versions.
In the version with the sample splitting, which is the one assumed in
Theorem~\ref{thm:5}, we use one half of the observations for the NR estimation and
the other half for the training. In the practical version, we use all the
observations for both steps.
\end{itemize}

We assumed that the number of the spikes is known and equal to one. We took
$d=20000$, $n_1=12$, $n_2=6$, $\delta=0.05$ and $100$ replications. We note that
the sample sizes are imbalanced, so that the bias correction is meaningful. The
implementation is given in \texttt{spiked\_svm\_classifiers.py}. We
give the results in Table~\ref{tab:compare}.

\begin{table}[ht]
\centering
\caption{The average of $e(1)+e(2)$ for $d=20{,}000$, $n_1=12$, $n_2=6$
and $\delta=0.05$. Standard errors in parentheses.}
\label{tab:compare}
\setlength{\tabcolsep}{3pt}
\begin{tabular}{@{}lrrrrr@{}}
\toprule
 & SVM & BC-SVM($\kappa$) & BC-SVM($\kappa_*$) & SC-SVM & SC-SVM (split) \\
\midrule
non-spiked
  & 1.0000 (0.0000) & 0.0000 (0.0000) & 0.0000 (0.0000)
  & 0.0015 (0.0001) & 0.0000 (0.0000) \\
$c=0.2$
  & 0.7887 (0.0135) & 0.2504 (0.0255) & 0.1600 (0.0200)
  & 0.1973 (0.0206) & 0.3684 (0.0279) \\
$c=0.5$
  & 0.3137 (0.0191) & 0.4071 (0.0383) & 0.1062 (0.0175)
  & 0.1185 (0.0171) & 0.2858 (0.0273) \\
$c=0.8$
  & 0.0347 (0.0099) & 0.4946 (0.0454) & 0.0220 (0.0082)
  & 0.0170 (0.0055) & 0.0710 (0.0131) \\
\bottomrule
\end{tabular}
\end{table}

We observe the following. In the non-spiked case, the SVM gives
$e(1)+e(2)=1.0000$, which is the strong inconsistency caused by the imbalance of
the sample sizes, and the bias correction removes it completely. This reproduces
the result of \citet{nakayama2017}. Under the spiked model, on the other hand,
BC-SVM($\kappa$) becomes worse as $c$ increases, and for $c=0.8$ it is much worse
than the SVM without any correction. This is the overcorrection given in
Proposition~\ref{prop:1}. In contrast, BC-SVM($\kappa_*$) gives a good performance
for all $c$.

In order to see the dependence on the sample size, we took $d=20000$, $c=0.8$ and
$\delta=0.01$, which is a harder setting, and varied $(n_1,n_2)$. We give the
results in Figure~\ref{fig:n-scale}.

\begin{figure}[ht]
\centering
\includegraphics[width=0.72\textwidth]{fig_n_scale.pdf}
\caption{Average of $e(1)+e(2)$ for $d=20{,}000$, $c=0.8$, $\delta=0.01$, as
$(n_1,n_2)$ grows.}
\label{fig:n-scale}
\end{figure}

We emphasize the behaviour of BC-SVM($\kappa$). Its misclassification rate is
$0.8759$, $0.8393$ and $0.7699$ at $(n_1,n_2)=(12,6)$, $(24,12)$ and $(48,24)$,
respectively, that is, it does not improve even though the sample size is
multiplied by four. This is because the overcorrection given in
Proposition~\ref{prop:1} is a systematic bias which does not vanish as the sample
size increases. On the other hand, BC-SVM($\kappa_*$) and the SC-SVM improve
rapidly. Hence Proposition~\ref{prop:1} is essential in practice.

We give the following remarks, which we consider to be informative for the users.

\begin{enumerate}[label=(\roman*)]
\item The SC-SVM and BC-SVM($\kappa_*$) give similar performances in
Table~\ref{tab:compare} and in Figure~\ref{fig:n-scale}, and the SC-SVM is better
only for $c=0.8$ in Table~\ref{tab:compare} and for the largest sample size in
Figure~\ref{fig:n-scale}. This is because the spike is orthogonal to $\mu$ in our
setting, so that the main damage caused by the spike comes through the bias term,
which BC-SVM($\kappa_*$) already removes. We expect that the advantage of the
SC-SVM is larger when the spike is not orthogonal to $\mu$ and the projection
removes a genuine source of noise in the discriminant direction.

\item The version with the sample splitting is worse than the practical version in
all the cases. This is because the splitting halves the number of the observations
used for the training, which is expensive when the sample size is small. We use
the splitting in Theorem~\ref{thm:5} in order to make the projected data independent
of the training data. We recommend the practical version in applications.

\item The gap between the two versions of the SC-SVM decreases as the sample size
increases, that is, it is $0.2603$, $0.0222$ and $0.0048$ at
$(n_1,n_2)=(12,6)$, $(24,12)$ and $(48,24)$, respectively. The gap is consistent
with Proposition~\ref{prop:3} in the sense that the estimation of the spiked
direction is poor when the sample size is small, so that the projection does not
work well.
\end{enumerate}
